% Source-derived chapter 03: Quot schemes and some of their basic properties
% Original source: rigid-local-systems-multiplicative-eigenvalue-problem
\chapter{Quot schemes and some of their basic properties}
\label{rigid-local-systems-multiplicative-eigenvalue-problem-chapter-03}
\source{rigid-local-systems-multiplicative-eigenvalue-problem}

\begin{remark}[Source section]
\label{rigid-local-systems-multiplicative-eigenvalue-problem-chapter-03-source}
\source{rigid-local-systems-multiplicative-eigenvalue-problem}
\source{rigid-local-systems-multiplicative-eigenvalue-problem}
This chapter reproduces the corresponding section of the retrieved
LaTeX source, including its definitions, statements, examples, and
proofs. It has not yet been translated into Lean.
\notready
\end{remark}

% The source section title was: Quot schemes and some of their basic properties
\label{quot1}
\source{rigid-local-systems-multiplicative-eigenvalue-problem}
We use notation and definitions  from the introduction.
\begin{defi}
\source{rigid-local-systems-multiplicative-eigenvalue-problem}
\notreadySuppose $\mn$ is a line bundle on $\Bbb{P}^1$. Let $\op{Quot}(d,r,\mn,n)\to \op{Bun}_{\mn}(n)$ denote the representable relative quot scheme of quotients, whose fiber over $\mw\in \op{Bun}_{\mn}(n)$ is the quot scheme of flat quotients $\mw\to\mq$ where $\mq$ is a coherent sheaf of rank $n-r$ and degree $d+\deg(\mn)$.
\end{defi}
See \cites{Laumon,CF} for the proof of the following:
\begin{proposition}
\source{rigid-local-systems-multiplicative-eigenvalue-problem}
\notready\label{laumon}
\source{rigid-local-systems-multiplicative-eigenvalue-problem}
$\Quot(d,r,\mn,n)$ is a smooth and irreducible Artin stack of dimension $r(n-r) +\deg(\mn)r +dn +1-n^2.$
\end{proposition}

\begin{defi}
\source{rigid-local-systems-multiplicative-eigenvalue-problem}
\notready\label{mindful}
\source{rigid-local-systems-multiplicative-eigenvalue-problem}
Let $\vec{J}=(J^1,\dots,J^s)$ be an $s$-tuple of subsets  of $[n]$ each of cardinality $r$.
\begin{enumerate}
\item Let $\Omega=\Omega(d,r,\mn,n,\vec{J})$ be the stack parametrizing
tuples $(\mv,\mw,\me,\gamma)$ where $\mw\to \mw/\mv$ gives a point of $\Quot(d,r,\mn,n)$, so that $\mv$ is a rank $r$ bundle of degree $-d$, and for each $p=p_i\in S$, the map
 $$\mv_p\to \mw_p/E^p_{j^i_k}$$
 has kernel of rank $\geq k$ for $k=1,\dots,r$.
\item Let $\Omega^0=\Omega^0(d,r,\mn,n,\vec{J})$ be the substack of $\Omega$ parametrizing tuples $(\mv,\mw,\me,\gamma)$ with $\mv$ a subbundle of $\mw$, and $\mv_p\in \Omega^0_{J^i}(\me_p)\subseteq \Gr(r,\mw_p)$ for all $p=p_i\in S$.
\end{enumerate}
\end{defi}
\begin{lemma}
\source{rigid-local-systems-multiplicative-eigenvalue-problem}
\notready\label{repres}
\source{rigid-local-systems-multiplicative-eigenvalue-problem}
The map $\Omega(d,r,\mn,n,\vec{J})\to \Par_{n,\mn,S}$, which sends $(\mv,\mw,\me,\gamma)$ to $(\mw,\me,\gamma)$, is representable and proper.
\end{lemma}
\begin{proof}
Let $\mathcal{Q}'$ be the fiber product of  $\Quot(d,r,\mn,n)$ and  $\Par_{n,\mn,S}$ over $\Bun_{\mathcal{N}}(n)$. The map  $\Omega(d,r,\mn,n,\vec{J})\to \mathcal{Q}'$ is representable and proper since the fibers are closed subschemes of products of flag varieties. Since  $\mathcal{Q}'\to \Par_{n,\mn,S}$ is a base change of $\Quot(d,r,\mn,n)\to \Bun_{\mathcal{N}}(n)$, it is proper and representable. The map $\Omega(d,r,\mn,n,\vec{J})\to \Par_{n,\mn,S}$ is a composite of these two maps and is hence  proper and representable.

\end{proof}

The following is the main result of this section. Statements of this form were first proved in \cite{bertie}:
 \begin{proposition}
\source{rigid-local-systems-multiplicative-eigenvalue-problem}
\notready\label{maya}
\source{rigid-local-systems-multiplicative-eigenvalue-problem}
  $\Omega=\Omega(d,r,\mn,n,\vec{J})$ is irreducible (possibly singular) of dimension
 \begin{equation}\label{codeyy}
\source{rigid-local-systems-multiplicative-eigenvalue-problem}
\dim \Quot(d,r,\mn,n) + \dim \Fl(n)^s -\sum_{i=1}^s |\sigma_{J^i}|
\end{equation}
Furthermore, $\Omega^0(d,r,\mn,n,\vec{J})$ is dense in $\Omega$.
 \end{proposition}
\subsection{Proof of Proposition \ref{maya}}
 By definition $\Omega$ is a subset of  tuples $(\mv,\mw,\me,\gamma)$ cut out by certain rank conditions on the maps $\mv_p\to \mw_p$ for $p=p_i\in S$. These rank conditions are given locally by $ |\sigma_{J^i}|$ equations for each $i$ (see the universal local case in the proof of \cite[Theorem 14.3]{FInt}), and hence the dimension of each irreducible component of $\Omega$ is at least as much as \eqref{codeyy}.

 It is also clear that $\Omega^0(d,r,\mn,n,\vec{J})$ is a smooth fiber bundle over an open subset of $\Quot(d,r,\mn,n)$ with fibres of dimension  $\dim \Fl(n)^s -\sum_{i=1}^s |\sigma_{J^i}|$ (Any fiber can be built by first considering induced flags on the  fibers of the subbundle and then extending these flags to the fibers of the universal bundle). Hence $\Omega^0$ has exactly the dimension \eqref{codeyy}.

 To conclude the proof we need to show that any irreducible component of $\Omega-\Omega^0$ has dimension less than \eqref{codeyy}.  We begin with  special cases in Sections \ref{bruck01}, \ref{bruck02} and \ref{bruck03}. The general case, which may be considered a linear superposition of these special cases,  is then proved in Section \ref{bruck04}. We also pay attention to  the divisorial components of  $\Omega-\Omega^0$ since they will play a special role in some computations in Section \ref{compact}.
 See \cite[Section 12.1]{BGM} for a similar computation.

\subsubsection{Subbundle property failing at points not in $S$}\label{bruck01}
\source{rigid-local-systems-multiplicative-eigenvalue-problem}
Consider first the following simple case: Let $q\not\in S$ and $\Omega_q(\epsilon)\subset \Omega$ the substack formed by
$(\mv,\mw,\me,\gamma)$ such that the following conditions are satisfied:
\begin{enumerate}
\item[(a)] $\mv\subseteq \mw$ is a subbundle on $\Bbb{P}^1-\{q\}$.
\iffalse
 This means that top exterior power of $\mv$ extends to a subbundle of an exterior power of $\mw$ after tensoring with a suitable $O(mq)$, and hence imposes algebraic conditions (so we have to include $m$ as well in the data).
 \fi
\item[(b)]   $\mv_p\in \Omega^0_{J^i}(\me_p)\subseteq \Gr(r,\mw_p)$ for all $p=p_i\in S$.
\item[(c)] The kernel of $\mv_q\to\mw_p$ has  dimension $\epsilon$.
\end{enumerate}
Let $\Quot_{\epsilon}(q)\subseteq  \Quot(d,r,\mn,n) $ be the substack formed by $\mv\subseteq \mw$ such that conditions (a) and (c) hold. Then the dimension of  $\Omega_q(\epsilon)$ is $\leq$ the quantity \eqref{codeyy} plus
\begin{equation}\label{simplecase}
\source{rigid-local-systems-multiplicative-eigenvalue-problem}
\dim \Quot_{\epsilon}(q)-\dim  \Quot(d,r,\mn,n)
\end{equation}

Consider a point $(\mv,\mw)$ in  $\Quot_{\epsilon}(q)$, and let $K$ be the kernel in (c). Let $\widetilde{\mv}$ be the subsheaf of $\mw$ which coincides with $\mv$ outside of $q$ and whose sections in a neighborhood of $q$ are meromorphic sections $s$ of $\mv$ such that $ts$ is a holomorphic section of $\mv$ with fiber at $q$ in $K$. Here $t$ is a uniformizing parameter at $q$. $\mv\subset\widetilde{\mv}\subseteq\mw$ and the image  $Q$ of $\mv_q\to \widetilde{\mv}_q$ has dimension $r-\epsilon$. The degree of $\widetilde{\mv}$ is $-d'$ where $d'=d-\epsilon$. The pairing $(\mv,\mw,K)$ to $(\widetilde{\mv},\mw,Q)$ is a bijection between triples with $K$ a subspace of the kernel of $\mv_q\to \mw_q$, and triples $(\widetilde{\mv},\mw,Q)$ with $Q$ a subspace of $\widetilde{\mv}_q$. The quantity
\eqref{simplecase} is therefore
$$\dim \Gr(r-\epsilon,r)+ (d'-d)n=\epsilon (r-n-\epsilon) =-(n+\epsilon-r)\epsilon\leq -2.$$
The point $q$ moves over a curve of dimension one. Therefore the following is a definite possibility:
\begin{itemize}
\item On a codimension one substack of $\Omega$, the quotient $\mw/\mv$ fails to be locally free. The only case where  this could happen is
$r=n-1$ and $\epsilon=1$.
\end{itemize}
\subsubsection{Subbundle property fails at a point of $S$}\label{bruck02}
\source{rigid-local-systems-multiplicative-eigenvalue-problem}
Consider first the following simple case: Let $p=p_1\in S$ and $\Omega_{p,L}(\delta)\subset \Omega$ the substack formed by
$(\mv,\mw,\me,\gamma)$ such that the following conditions hold:
\begin{enumerate}
\item[(a)] $\mv\subseteq \mw$ is a subbundle on $\Bbb{P}^1-\{p\}$.
\item[(b)]   $\mv_{p'}\in \Omega^0_{J^i}(\me_{p'})\subseteq \Gr(r,\mw_{p'})$ for all $p'\neq p\in S$.
\item[(c)] $\mv_q\to\mw_q$ has a kernel of dimension $\delta$.
\item[(d)] The image $Q$ of $\mv_p\to\mw_p$ is an $r-\delta$ dimensional subspace of $\mw_p$ which lies in the Schubert cell  corresponding to $L=\{l_1,\dots,l_{r-\delta}\}$.
\end{enumerate}

Computing as before, the dimension of  $\Omega_{p,L}(\delta)$ is $\leq$ the quantity \eqref{codeyy} plus
\begin{equation}\label{simplecase1}
\source{rigid-local-systems-multiplicative-eigenvalue-problem}
\dim \Gr(r-\delta,r)+ (d'-d)n + |\sigma_L|-|\sigma_{J^1}|
\end{equation}
We also use $l_a\leq j^1_{a+\delta}$ for $a\leq r-\delta$ (this follows from the rank conditions).  This gives  $|\sigma_L|-|\sigma_{J^1}|\leq -\sum_{a=1}^{\delta} (n-r+a-j^1_a)$. Therefore the quantity \eqref{simplecase1} is $\leq$
$\sum_{a=1}^{\delta} (a-j^1_a-\delta)$. This number is $\leq -2$ except for when $\delta=1$ and $j^1_1=1$ when it is equal to $-1$.

\subsubsection{Schubert condition strengthens at a point of $S$}\label{bruck03}
\source{rigid-local-systems-multiplicative-eigenvalue-problem}
We need to now consider the cases in which $\mv$ is a subbundle of $\mw$, but the Schubert condition at $p=p_1$ strengthens to a Schubert cell parametrized by a subset $J$ of $[n]$, $|J|=r$. Dimension counts still apply, and we see that the codimension drops by $1$ whenever $J$ is of the form $J=(I^1-\{a\})\cup \{a-1\}$ where $a>1$, $a\in I^1$, $a-1\not\in I^1$.
\subsubsection{Conclusion of the proof of Proposition \ref{maya}}\label{bruck04}
\source{rigid-local-systems-multiplicative-eigenvalue-problem}
It remains to  show that any irreducible component  of $\Omega-\Omega^0$ has dimension less than the quantity \eqref{codeyy}.  Consider the generic point $(\mv,\mw,\me,\gamma)$ of an irreducible component $Y$.
\begin{enumerate}
\item Let $Q=\{q_1,\dots,q_m\}$ be the set of points (possibly empty) not in $S$ such that for $q\in Q$, $\mv_q\to \mw_q$ has kernel of rank $\epsilon(q)>0$.
\item Let $\delta: S\to \Bbb{Z}_{\geq 0}$ be such that the maps $\mv_p\to\mw_p$ have kernels $K(p)$ of rank $\delta(p)$, for $p\in S$. Assume further that the images $Q(p)$ of $\mv_p\to\mw_p$ are  $r-\delta(p)$ dimensional subspaces of $\mw_q$  in the Schubert cells corresponding to $L^p=\{l^p_1,\dots,l^p_{r-\delta(p)}\}$.
\end{enumerate}
Computing as before, the dimension of  $Y$ is $\leq$ the quantity \eqref{codeyy} plus
$$\sum_{q\in Q} \bigl(1-(n+\epsilon(q)-r)\epsilon(q)\bigr) + \sum_{p\in S} \sum_{a=1}^{\delta(p)} (a-j^p_a-\delta(p)).$$
Each of the $q$ contributions is $\leq -1$ by the computation in Section \ref{bruck01}. If $\delta(p)>0$, then the $p$ term is $\leq -1$ by the computation in Section \ref{bruck02}. Note that if there no $q$ terms and $\delta(p)=0$ for all $p$, then since $Y\cap\Omega^0=\emptyset$, we need $L^p\neq J^p$ for some $p$, and this makes the dimension of $Y$ strictly less than  the quantity \eqref{codeyy} as in Section \ref{bruck03}. This corresponds to one of the inequalities $l^p_a\leq j^p_{a+\delta(p)}$ in Section \ref{bruck02} becoming strict.

Therefore the dimension of $Y$ is strictly less than the quantity \eqref{codeyy}, and this completes the proof of Proposition \ref{maya}.
\subsection{Cycle classes}
Propositions \ref{laumon}, \ref{maya} and Lemma \ref{repres} give the following
\begin{corollary}
\source{rigid-local-systems-multiplicative-eigenvalue-problem}
\notready
The representable morphism $\Omega(d,r,\mn,n,\vec{J})\to \Par_{n,\mn,S}$ has relative dimension
$\sum_{i=1}^s |\sigma_{J^i}| -(dn +\deg(\mn)r +r(n-r)).$
\end{corollary}
For proper morphisms of irreducible schemes $Y\to X$ of relative dimension $d$, one can form the push forward of the fundamental class $[Y]$  of $Y$, giving  a codimension $d$ cycle on $X$ \cite[Section 1.4]{FInt}. The push forward operation commutes with flat base changes by \cite[Proposition 1.7]{FInt}, and is therefore applicable in the setting of representable morphisms of irreducible Artin stacks. Applying it to the morphism $\Omega(d,r,\mn,n,\vec{J})\to \Par_{n,\mn,S}$, we obtain a cycle class on $\Par_{n,\mn,S}$:
\begin{defi}
\source{rigid-local-systems-multiplicative-eigenvalue-problem}
\notready\label{cyclist}
\source{rigid-local-systems-multiplicative-eigenvalue-problem}
The pushforward of the cycle $\Omega(d,r,\mn,n,\vec{J})$ to $\Par_{n,\mn,S}$ gives an effective cycle $C(d,r,\mn,n,\vec{J})$ (possibly zero) of codimension
 $\sum_{i=1}^s |\sigma_{J^i}| -(dn +\deg(\mn)r +r(n-r)).$
 \end{defi}
